\section{Математическое описание}

В этом разделе описана модель автомата: поле, состояния клеток, правило их
изменения и способ перехода от одного шага к следующему.

\subsection{Клеточный автомат}

Клеточные автоматы – это регулярная структура динамических объектов (клеток),
функционирующих синхронно. Клеточные автоматы моделируют процессы,
разворачивающиеся в дискретном пространстве и в дискретном времени

Клеточный автомат имеет состояние, определяемое как вектор или матрица
состояний компонентных автоматов. Каждый автомат функционирует как автомат Мура,
т.е. это автомат без выхода, выходом является его состояние. Вход на каждом шаге
клеточного автомата – состояния его соседей. На следующем шаге (в следующем такте)
новое состояние каждого автомата определяется как функция его текущего состояния и
текущих состояний его соседей


Наш автомат можно описать набором $A = (L, S, N, f)$, где:

\begin{itemize}
    \item $L$ --- квадратная сетка размером $N \times N$;
    \item $S = \{0, 1\}$ --- возможные состояния клетки: $0$ означает, что
    клетка мёртвая, а $1$ --- что живая;
    \item $N$ --- набор клеток, состояние которых учитывается при обновлении:
    сама клетка и четыре соседа сверху, справа, снизу и слева;
    \item $f$ --- правило, которое по состояниям этих пяти клеток определяет
    новое состояние центральной клетки.
\end{itemize}

Размер сетки обозначим $N$. Состояние клетки в строке $i$ и столбце $j$ на
шаге $t$ обозначим $s^t_{i,j}$. Индексы $i$ и $j$ принимают значения от $0$
до $N-1$. Начальное поле задаёт состояние автомата на шаге $t=0$.

\subsection{Правило перехода}

Состояние клетки на следующем шаге определяется её текущим состоянием и
состояниями четырёх соседей:

\[
s^{t+1}_{i,j} =
f\left(s^t_{i,j}, s^t_{i-1,j}, s^t_{i,j+1},
s^t_{i+1,j}, s^t_{i,j-1}\right).
\]

Здесь $f$ --- правило перехода, которое возвращает $0$ или $1$. Входные
значения в формуле идут в порядке: (центр, верх, право, низ, лево).

Пять клеток могут иметь $2^5=32$ разных набора состояний. Поэтому правило
хранится как число из 32 бит: каждому набору соответствует один бит, который
задаёт новое состояние центральной клетки. Для варианта $v$ число правила
вычисляется по формуле

\[
R = 2 \cdot 11 \cdot 2006 \cdot 13 \cdot 12. = 6\,884\,592_{10}
\]
\[
= 00000000011010010000110011110000_2;
\]
 Для набора
$(a,b,c,d,e)$, где значения идут в указанном выше порядке, его номер равен
$k = 16a + 8b + 4c + 2d + e$. Программа проверяет соответствующий этому
номеру бит числа $R$: если он равен 1, центральная клетка становится живой,
если 0 --- мёртвой.

Таблица истинности функции $f$ приведена в таблице~\ref{tab:truth-table}.
Столбцы идут в том же порядке, что и входы функции: центр, верх, право,
низ и лево.

\begin{longtable}{|c|c|c|c|c|c|}
\caption{Функция $6\,884\,592$}\label{tab:truth-table} \\
\hline
\makecell{$s_{i,j}$} &
\makecell{$s_{i-1,j}$} &
\makecell{$s_{i,j+1}$} &
\makecell{$s_{i+1,j}$} &
\makecell{$s_{i,j-1}$} &
\makecell{$f(s_{i,j}, s_{i-1,j}, s_{i,j+1},$\\
$s_{i+1,j}, s_{i,j-1})$} \\
\hline
\endfirsthead
\hline
\endhead
0 & 0 & 0 & 0 & 0 & 0 \\
0 & 0 & 0 & 0 & 1 & 0 \\
0 & 0 & 0 & 1 & 0 & 0 \\
0 & 0 & 0 & 1 & 1 & 0 \\
0 & 0 & 1 & 0 & 0 & 0 \\
0 & 0 & 1 & 0 & 1 & 0 \\
0 & 0 & 1 & 1 & 0 & 0 \\
0 & 0 & 1 & 1 & 1 & 0 \\
0 & 1 & 0 & 0 & 0 & 0 \\
0 & 1 & 0 & 0 & 1 & 1 \\
0 & 1 & 0 & 1 & 0 & 1 \\
0 & 1 & 0 & 1 & 1 & 0 \\
0 & 1 & 1 & 0 & 0 & 1 \\
0 & 1 & 1 & 0 & 1 & 0 \\
0 & 1 & 1 & 1 & 0 & 0 \\
0 & 1 & 1 & 1 & 1 & 1 \\
1 & 0 & 0 & 0 & 0 & 0 \\
1 & 0 & 0 & 0 & 1 & 0 \\
1 & 0 & 0 & 1 & 0 & 0 \\
1 & 0 & 0 & 1 & 1 & 0 \\
1 & 0 & 1 & 0 & 0 & 1 \\
1 & 0 & 1 & 0 & 1 & 1 \\
1 & 0 & 1 & 1 & 0 & 0 \\
1 & 0 & 1 & 1 & 1 & 0 \\
1 & 1 & 0 & 0 & 0 & 1 \\
1 & 1 & 0 & 0 & 1 & 1 \\
1 & 1 & 0 & 1 & 0 & 1 \\
1 & 1 & 0 & 1 & 1 & 1 \\
1 & 1 & 1 & 0 & 0 & 0 \\
1 & 1 & 1 & 0 & 1 & 0 \\
1 & 1 & 1 & 1 & 0 & 0 \\
1 & 1 & 1 & 1 & 1 & 0 \\
\hline
\end{longtable}

\subsection{Граничные условия}

В поле используются тороидальные граничные условия: противоположные края
считаются соединёнными. Поэтому сосед клетки у верхнего края находится в
нижней строке, у нижнего края --- в верхней строке. Аналогично соединены
левый и правый края.

При обращении к клетке за границей поля её координата возвращается внутрь
поля по остатку от деления на $N$. Например, строка $-1$ соответствует
строке $N-1$, а строка $N$ --- строке $0$.

В программе для этого используется функция \texttt{wrap}. Для значения value она вычисляет индекс:

\[
\operatorname{wrap}(value) = ((value \bmod N) + N) \bmod N.
\]

Здесь $N$ --- размер поля. Формула возвращает индекс от $0$ до $N-1$ даже
для отрицательных координат. Например,
$\operatorname{wrap}(-1) = N-1$.

\subsection{Смена поколений}

Чтобы получить новое поколение, правило применяется к каждой клетке поля.
Все новые состояния сначала вычисляются по старому полю и сохраняются
отдельно. После этого старое поле целиком заменяется новым. Поэтому при
расчёте шага изменение одной клетки не влияет на вычисление соседних клеток
в том же шаге.
